\subsection{4.1 分模块分析类图}
\noindent 分析类由用例名词、字段和事件抽取；先核对账号/人员泛化，再标组合、关联与基数。每张图只表达一个模块，跨组类以稳定引用相连。图中 `CartLine` 等短暂对象不自动等于持久表；目标分析类与当前 `users.role` 等实现差异在数据设计章注明。
% C-CUST 顾客与预约：分析类，不等同现有数据库表。
\newcommand{\BSRClassCustomer}{\begin{tikzpicture}[x=1cm,y=1cm]
  \BSRClassCard{account}{0}{0}{Account《抽象》}{id；role；status}{authenticate()；authorize()}
  \BSRClassCard{customer}{5.1}{0}{Customer}{customerId；email；量体}{updateProfile()；identify()}
  \BSRClassCard{reservation}{10.2}{0}{Reservation}{id；租期；状态；技能自评}{hold()；cancel()}
  \BSRClassCard{cart}{0}{-4.2}{CartLine}{规格；数量；报价版本}{changeQty()；invalidateQuote()}
  \BSRClassCard{item}{5.1}{-4.2}{ReservationItem}{规格；数量；单价快照}{freezeQuote()；reprice()}
  \BSRClassCard{consent}{10.2}{-4.2}{PolicyConsent}{政策版本；同意时间}{recordConsent()；verify()}
  \draw[Navy,line width=.7pt,-{Triangle[open]}] (customer) -- node[above,font=\tiny,fill=white]{泛化 1:1} (account);
  \draw[Navy,line width=.7pt] (customer) -- node[above,font=\tiny,fill=white]{1:n 预约} (reservation);
  \draw[Navy,line width=.7pt,{Diamond[fill=Navy,length=2.5mm,width=2.2mm]}-] (reservation) -- node[sloped,above,font=\tiny,fill=white]{组合 1:n} (item);
  \draw[Navy,line width=.7pt] (reservation) -- node[right,font=\tiny,fill=white]{1:1 同意} (consent);
  \draw[Navy,dashed,line width=.7pt,-{Latex}] (cart) -- node[above,font=\tiny,fill=white]{提交后生成} (item);
  \node[bsrNote] at (5.1,-7.2) {来源：UC-01～04。CartLine 可为短暂对象；ReservationItem 与 PolicyConsent 才保留历史承诺。};
\end{tikzpicture}}

% C-BOOK 商品规格、物理单件与门店。
\newcommand{\BSRClassInventory}{\begin{tikzpicture}[x=1cm,y=1cm]
  \BSRClassCard{category}{0}{0}{Category}{id；name；parentId}{addChild()；checkCycle()}
  \BSRClassCard{variant}{5.1}{0}{Variant}{id；款式；尺码；技能}{matches()；quote()}
  \BSRClassCard{unit}{10.2}{0}{Unit}{unitId；库存号；状态}{isAvailable()；isolate()}
  \BSRClassCard{store}{0}{-4.2}{Store}{storeId；名称；营业时段}{receive()；handover()}
  \BSRClassCard{location}{5.1}{-4.2}{UnitLocation}{当前店；归属店；版本}{moveCurrent()；transferHome()}
  \BSRClassCard{movement}{10.2}{-4.2}{InventoryMovement}{unitId；来源/目标店；时间}{append()；audit()}
  \BSRClassCard{price}{5.1}{-8.4}{PricePolicy}{规格/门店；日价 NOK；版本}{effectiveAt()；snapshot()}
  \draw[Navy,line width=.7pt] (category) -- node[above,font=\tiny,fill=white]{1:n 分类} (variant);
  \draw[Navy,line width=.7pt] (variant) -- node[above,font=\tiny,fill=white]{1:n 实物} (unit);
  \draw[Navy,line width=.7pt] (category.west) to[out=165,in=115,looseness=2] node[left,font=\tiny]{自反 0:n} (category.north);
  \draw[Navy,line width=.7pt] (store) -- node[above,font=\tiny,fill=white]{1:n 位置} (location);
  \draw[Navy,line width=.7pt] (unit) -- node[sloped,above,font=\tiny,fill=white]{1:1 位置} (location);
  \draw[Navy,line width=.7pt] (unit) -- node[right,font=\tiny,fill=white]{1:n 流水} (movement);
  % 价格政策从规格外侧绕行，避免纵线穿过 UnitLocation 类框。
  \coordinate (priceRouteTop) at ([xshift=7mm]variant.east);
  \coordinate (priceRouteBottom) at (priceRouteTop |- price.east);
  \draw[Navy,line width=.7pt] (variant.east) -- (priceRouteTop) -- (priceRouteBottom) -- (price.east);
  \node[font=\tiny,rotate=90,fill=white,inner sep=1pt] at ($(priceRouteTop)!0.73!(priceRouteBottom)$) {1:n 价目版本};
  \node[bsrNote] at (10.2,-8.4) {来源：UC-02/07/09/13。规格用于筛选，单件用于交接；异店归还只改当前店。};
\end{tikzpicture}}

% C-CONT 合同版本、备货和逐件履约。
\newcommand{\BSRClassContract}{\begin{tikzpicture}[x=1cm,y=1cm]
  \BSRClassCard{contract}{0}{0}{Contract}{id；合同号；版本；签署状态}{issue()；signCurrent()}
  \BSRClassCard{line}{5.1}{0}{ContractLine}{规格；数量；NOK 快照}{freeze()；verifyTotal()}
  \BSRClassCard{adjust}{10.2}{0}{ContractAdjustment}{原/新版；差额；同意}{preview()；commit()}
  \BSRClassCard{prep}{0}{-4.2}{Preparation}{reservationId；指派人；状态}{assign()；markReady()}
  \BSRClassCard{prepu}{5.1}{-4.2}{PreparationUnit}{订单行；unitId；检查}{bind()；safetyCheck()}
  \BSRClassCard{rentalu}{10.2}{-4.2}{RentalUnit}{订单行；unitId；借/还时刻}{handover()；receive()}
  \draw[Navy,line width=.7pt,{Diamond[fill=Navy,length=2.5mm,width=2.2mm]}-] (contract) -- node[above,font=\tiny,fill=white]{组合 1:n} (line);
  % 合同至调整版本的关联走上方，不穿过中间的 ContractLine。
  \coordinate (adjustRouteL) at ([yshift=7mm]contract.north);
  \coordinate (adjustRouteR) at ([yshift=7mm]adjust.north);
  \draw[Navy,line width=.7pt] (contract.north) -- (adjustRouteL) -- node[midway,above,font=\tiny,fill=white]{1:n 版本调整} (adjustRouteR) -- (adjust.north);
  \draw[Navy,line width=.7pt,{Diamond[fill=Navy,length=2.5mm,width=2.2mm]}-] (prep) -- node[above,font=\tiny,fill=white]{组合 1:n} (prepu);
  \draw[Navy,line width=.7pt] (line) -- node[right,font=\tiny,fill=white]{1:n 单件} (prepu);
  \draw[Navy,line width=.7pt] (prepu) -- node[above,font=\tiny,fill=white]{1:1 交付} (rentalu);
  \node[bsrNote] at (5.1,-7.2) {来源：UC-05～09。预约行与实物为时间化 m:n；当前签署版本是交付守卫。};
\end{tikzpicture}}

% C-SERV 归还事实、受损鉴定、维修与安全验收。
\newcommand{\BSRClassDamage}{\begin{tikzpicture}[x=1cm,y=1cm]
  \BSRClassCard{inspection}{0}{0}{Inspection}{unitId；事实分类；实还时刻}{recordFact()；isolate()}
  \BSRClassCard{damage}{5.1}{0}{DamageCase}{等级；状态；24h 截止}{assess()；requestReview()}
  \BSRClassCard{photo}{10.2}{0}{DamagePhoto}{部位；受控图片；上传人}{attach()；authorizeRead()}
  \BSRClassCard{repair}{0}{-4.2}{MaintenanceOrder}{unitId；技师；费率快照}{diagnose()；submit()}
  \BSRClassCard{rate}{5.1}{-4.2}{ContractorRate}{技师；生效日；小时 NOK}{selectAt()；snapshot()}
  \BSRClassCard{safety}{10.2}{-4.2}{SafetyAcceptance}{工单；检查项；经理结论}{approve()；returnForFix()}
  \draw[Navy,line width=.7pt] (inspection) -- node[above,font=\tiny,fill=white]{0:1 疑损} (damage);
  \draw[Navy,line width=.7pt,{Diamond[fill=Navy,length=2.5mm,width=2.2mm]}-] (damage) -- node[above,font=\tiny,fill=white]{组合 1:n} (photo);
  \draw[Navy,line width=.7pt] (damage) -- node[sloped,above,font=\tiny,fill=white]{0:1 维修} (repair);
  \draw[Navy,line width=.7pt] (rate) -- node[above,font=\tiny,fill=white]{历史费率} (repair);
  % 维修至安全验收的关联走下方，不穿过 ContractorRate。
  \coordinate (safetyRouteL) at ([yshift=-6mm]repair.south);
  \coordinate (safetyRouteR) at ([yshift=-6mm]safety.south);
  \draw[Navy,line width=.7pt] (repair.south) -- (safetyRouteL) -- node[midway,above,font=\tiny,fill=white]{1:n 验收轮次} (safetyRouteR) -- (safety.south);
  \node[bsrNote] at (5.1,-7.2) {来源：UC-09/10/12。专员记事实，技师初评，经理核费与安全；“修好”不等于“可租”。};
\end{tikzpicture}}

% C-FIN 费用、押金与资金流水；正式金额只有 NOK。
\newcommand{\BSRClassFinance}{\begin{tikzpicture}[x=1cm,y=1cm]
  \BSRClassCard{late}{0}{0}{LateCharge}{unitId；实还时间；规则版本}{estimate()；freeze()}
  \BSRClassCard{damagec}{5.1}{0}{DamageCharge}{鉴定号；核定金额；依据}{review()；approve()}
  \BSRClassCard{settle}{10.2}{0}{DepositSettlement}{押金；抵扣；可退；补差}{calculate()；holdDispute()}
  \BSRClassCard{pay}{0}{-4.2}{PaymentEvent}{类型；金额 NOK；回执状态}{recordOnce()；verify()}
  \BSRClassCard{contact}{5.1}{-4.2}{SupplementContact}{责任专员；联系结果；时间}{logContact()；close()}
  \BSRClassCard{rule}{10.2}{-4.2}{PolicySnapshot}{超时/损坏版；同意时间}{resolve()；preserve()}
  % 超时费用至押金结算的关联走上方，不穿过 DamageCharge。
  \coordinate (settleRouteL) at ([yshift=7mm]late.north);
  \coordinate (settleRouteR) at ([yshift=7mm]settle.north);
  \draw[Navy,line width=.7pt] (late.north) -- (settleRouteL) -- node[midway,above,font=\tiny,fill=white]{0:n 汇总} (settleRouteR) -- (settle.north);
  \draw[Navy,line width=.7pt] (damagec) -- node[above,font=\tiny,fill=white]{0:n 汇总} (settle);
  \draw[Navy,line width=.7pt] (settle) -- node[sloped,above,font=\tiny,fill=white]{1:n 资金事件} (pay);
  \draw[Navy,line width=.7pt] (settle) -- node[above,font=\tiny,fill=white]{0:n 补差联系} (contact);
  \draw[Navy,line width=.7pt] (rule) -- node[right,font=\tiny,fill=white]{1:1 冻结规则} (settle);
  \node[bsrNote] at (5.1,-7.2) {来源：UC-03/05/08/11。费用扣抵按整数 NOK；争议暂停结清，模拟回执不等于真实到账。};
\end{tikzpicture}}

% C-STAFF 人员/排班/提醒。账号为共同身份，技师有对象级任务授权。
\newcommand{\BSRClassWorkforce}{\begin{tikzpicture}[x=1cm,y=1cm]
  \BSRClassCard{identity}{0}{0}{WorkerIdentity《抽象》}{accountId；角色；状态}{authenticate()；authorize()}
  \BSRClassCard{employee}{5.1}{0}{Employee}{employeeId；所属店；资料}{viewShift()；clock()}
  \BSRClassCard{technician}{10.2}{0}{Technician}{contractorId；专业；费率}{acceptTask()；recordWork()}
  \BSRClassCard{shift}{0}{-4.2}{Shift}{日期；工作店；起止}{assign()；rejectConflict()}
  \BSRClassCard{punch}{5.1}{-4.2}{Punch}{上/下班时刻；工作日}{clockIn()；clockOut()}
  \BSRClassCard{alert}{10.2}{-4.2}{OrderAlert}{对象；责任人；已读/未结}{markRead()；resolveCase()}
  \BSRClassCard{audit}{5.1}{-8.4}{AuditEvent}{操作者；对象；前后状态}{append()；query()}
  \draw[Navy,line width=.7pt,-{Triangle[open]}] (employee) -- node[above,font=\tiny,fill=white]{泛化} (identity);
  % 技师泛化线绕过 Employee 类框，空心三角仍指向抽象身份类。
  \coordinate (techRouteR) at ([yshift=7mm]technician.north);
  \coordinate (techRouteL) at ([yshift=7mm]identity.north);
  \draw[Navy,line width=.7pt,-{Triangle[open]}] (technician.north) -- (techRouteR) -- node[midway,above,font=\tiny,fill=white]{泛化} (techRouteL) -- (identity.north);
  \draw[Navy,line width=.7pt] (employee) -- node[sloped,above,font=\tiny,fill=white]{1:n 计划} (shift);
  \draw[Navy,line width=.7pt] (employee) -- node[right,font=\tiny,fill=white]{1:n 事实} (punch);
  \draw[Navy,line width=.7pt] (identity) -- node[sloped,above,font=\tiny,fill=white]{0:n 授权提醒} (alert);
  \draw[Navy,line width=.7pt] (identity) -- node[pos=.82,right,font=\tiny,fill=white]{1:n 审计} (audit);
  \node[bsrNote] at (10.2,-8.4) {来源：UC-14/15。班次是计划，打卡是事实；已读提醒与订单未解决分离。};
\end{tikzpicture}}

\BSRNewFigure{\BSRClassCustomer}{C-CUST 顾客与预约分析类：UC-01～04，账号泛化、预约行与同意版本}{fig:c-cust}
\BSRNewFigure{\BSRClassInventory}{C-BOOK 库存与门店分析类：UC-02/07/09/13，规格与物理单件分离}{fig:c-book}
\BSRNewFigure{\BSRClassContract}{C-CONT 合同与逐件履约分析类：UC-05～09，版本与时间化连接}{fig:c-cont}
\BSRNewFigure{\BSRClassDamage}{C-SERV 归还事实、受损鉴定与维修分析类：UC-09/10/12}{fig:c-serv}
\BSRNewFigure{\BSRClassFinance}{C-FIN 费用、押金和资金事件分析类：UC-03/05/08/11}{fig:c-fin}
\BSRNewFigure{\BSRClassWorkforce}{C-STAFF 人员、排班、提醒与审计分析类：UC-14/15}{fig:c-staff}
\noindent C-MINI 仅增加微信身份绑定、一次性挑战和独立会话；顾客账号、购物车与订单仍引用已有 C-CUST/C-CONT 聚合。其余角色没有小程序工作台。
\BSRDiagram{019-class-identity}{C-MINI 微信身份分析类：边界、控制、绑定与会话分离；UC-17/21}{fig:c-mini}

\begin{BSRTable}
\rowcolor{TableHead}\color{white}\bfseries 关系检查 & \color{white}\bfseries 类与用例来源、当前/目标差异 \\
顾客/员工/技师泛化 & 共同身份、认证和审计抽为 `Account` / `WorkerIdentity`；目标分析可画子类，当前演示以 `users.role` 单表表达，不把子类当已建表。\\
预约行与单件 & UC-02/04 的预约行先记规格和数量；UC-07/09 的准备单件、交接单件是时间化连接，允许同一单件在不重叠租期服务多张订单。\\
合同与政策快照 & UC-05/11 的合同版本、NOK 分项、规则同意和技师费率要冻结旧值，后来改价不回写历史。\\
受损/维修责任 & UC-10/12 中专员记录检查事实，技师初评和维修，经理复核费用与安全验收；技师点击提交不能直接使单件可租。\\
人员计划与事实 & UC-14 的 `Shift` 是计划、`Punch` 是服务器记录的实际打卡；UC-15 的提醒已读与订单未解决是两个字段。\\
\end{BSRTable}
